% =============================================================================
% Figure blocks for Mishal's diagrams -- copy each block into the chapter where
% it belongs (suggested: use-case -> ch4; domain model, ERD, anchor -> ch5).
%
% Preamble of main.tex needs (once):
%     \usepackage{graphicx}
%     \usepackage{pdflscape}
%     % =============================================================================
% UNIFY -- shared TikZ styles for the SE diagrams
% Load ONCE in the report preamble (main.tex), after \usepackage{tikz-uml}:
%     % =============================================================================
% UNIFY -- shared TikZ styles for the SE diagrams
% Load ONCE in the report preamble (main.tex), after \usepackage{tikz-uml}:
%     % =============================================================================
% UNIFY -- shared TikZ styles for the SE diagrams
% Load ONCE in the report preamble (main.tex), after \usepackage{tikz-uml}:
%     \input{diagrams/unify-diagram-styles}
% Plain TikZ only (no tikz-uml dependency), so it compiles on any TeX Live.
% =============================================================================
\usepackage{tikz}
\usetikzlibrary{arrows.meta,positioning,calc,fit,backgrounds,
                shapes.geometric,shapes.multipart,decorations.pathreplacing}

% ---- palette (prints legibly in greyscale) ----------------------------------
\definecolor{uInk}{HTML}{1F2933}
\definecolor{uMuted}{HTML}{5B6670}
\definecolor{uLine}{HTML}{3E4C59}
\definecolor{uIngest}{HTML}{D9E8F5}   % event reliability / version store
\definecolor{uGraph}{HTML}{DDEFDF}    % evidence graph / policy / traceability
\definecolor{uVerify}{HTML}{F6E6CF}   % verification / findings
\definecolor{uIntent}{HTML}{E7E1F3}   % intent-side artifacts
\definecolor{uOk}{HTML}{2F7D4E}
\definecolor{uWarn}{HTML}{B7791F}
\definecolor{uBad}{HTML}{B4372F}
\definecolor{uSpan}{HTML}{FCE9A6}

% ---- fixed font sizes, so layouts don't break in 11pt/12pt report classes ----
\newcommand{\uDiagScript}{\fontsize{7pt}{8.5pt}\selectfont}
\newcommand{\uDiagSmall}{\fontsize{9pt}{11pt}\selectfont}
\newcommand{\uDiagNormal}{\fontsize{10pt}{12pt}\selectfont}

\tikzset{
  % ---- generic ---------------------------------------------------------------
  every picture/.append style={line join=round,
     execute at begin picture={\let\scriptsize\uDiagScript
       \let\small\uDiagSmall\let\normalsize\uDiagNormal\uDiagSmall}},
  every node/.append style={execute at begin node={\hyphenpenalty=10000\exhyphenpenalty=10000}},
  uline/.style={draw=uLine, line width=0.6pt},
  lbl/.style={font=\scriptsize, text=uMuted, fill=white, inner sep=1pt},
  note/.style={draw=uMuted, fill=yellow!8, font=\scriptsize, align=left,
               text width=#1, inner sep=4pt, rounded corners=1pt},
  note/.default=4.6cm,
  % ---- use-case --------------------------------------------------------------
  usecase/.style={ellipse, draw=uLine, fill=white, line width=0.6pt,
                  minimum width=4.4cm, minimum height=0.95cm, align=center,
                  font=\scriptsize, inner sep=1pt, text width=3.6cm},
  sysboundary/.style={draw=uLine, line width=0.9pt, rounded corners=3pt,
                      fill=uIngest!25},
  assoc/.style={draw=uLine, line width=0.5pt},
  stereo/.style={draw=uLine, line width=0.5pt, dashed,
                 -{Stealth[length=2mm]}},
  extsys/.style={draw=uLine, fill=uGraph!60, rounded corners=2pt,
                 minimum width=2.3cm, minimum height=0.8cm, font=\scriptsize,
                 align=center},
  % ---- ERD / class boxes -------------------------------------------------------
  table/.style={rectangle split, rectangle split parts=2, draw=uLine,
                line width=0.6pt, rectangle split part fill={#1,white},
                text width=4.9cm, align=left, inner sep=3pt,
                font=\scriptsize\ttfamily},
  table/.default=uGraph,
  klass/.style={rectangle split, rectangle split parts=2, draw=uLine,
                line width=0.6pt, rectangle split part fill={#1,white},
                text width=3.7cm, align=center, inner sep=3pt,
                font=\scriptsize},
  klass/.default=white,
  rel/.style={draw=uLine, line width=0.55pt},
  card/.style={font=\scriptsize\bfseries, text=uInk, inner sep=1.5pt},
  generalise/.style={draw=uLine, line width=0.6pt,
                     -{Triangle[open, length=3mm, width=3mm]}},
  directed/.style={draw=uLine, line width=0.55pt, -{Stealth[length=2.2mm]}},
  % ---- flow ----------------------------------------------------------------------
  decision/.style={diamond, aspect=1.9, draw=uLine, fill=white,
                   line width=0.6pt, font=\scriptsize, align=center,
                   inner sep=1pt, text width=2.1cm},
  step/.style={rectangle, rounded corners=2pt, draw=uLine, line width=0.6pt,
               font=\scriptsize, align=center, minimum height=0.8cm,
               text width=#1, inner sep=3pt},
  step/.default=2.6cm,
  flow/.style={draw=uLine, line width=0.6pt, -{Stealth[length=2.2mm]}},
}

% ---- stick-figure actor:  \actor{name}{(x,y)}{Label} ---------------------------
\newcommand{\actor}[3]{%
  \begin{scope}[shift={#2}, line width=0.8pt, draw=uInk]
    \draw (0,0.62) circle (0.17);
    \draw (0,0.45) -- (0,-0.05);
    \draw (-0.32,0.28) -- (0.32,0.28);
    \draw (0,-0.05) -- (-0.24,-0.45);
    \draw (0,-0.05) -- (0.24,-0.45);
    \node[font=\scriptsize\bfseries, align=center, anchor=north]
         (#1-label) at (0,-0.5) {#3};
    \coordinate (#1) at (0.4,0.2);
  \end{scope}}

% Plain TikZ only (no tikz-uml dependency), so it compiles on any TeX Live.
% =============================================================================
\usepackage{tikz}
\usetikzlibrary{arrows.meta,positioning,calc,fit,backgrounds,
                shapes.geometric,shapes.multipart,decorations.pathreplacing}

% ---- palette (prints legibly in greyscale) ----------------------------------
\definecolor{uInk}{HTML}{1F2933}
\definecolor{uMuted}{HTML}{5B6670}
\definecolor{uLine}{HTML}{3E4C59}
\definecolor{uIngest}{HTML}{D9E8F5}   % event reliability / version store
\definecolor{uGraph}{HTML}{DDEFDF}    % evidence graph / policy / traceability
\definecolor{uVerify}{HTML}{F6E6CF}   % verification / findings
\definecolor{uIntent}{HTML}{E7E1F3}   % intent-side artifacts
\definecolor{uOk}{HTML}{2F7D4E}
\definecolor{uWarn}{HTML}{B7791F}
\definecolor{uBad}{HTML}{B4372F}
\definecolor{uSpan}{HTML}{FCE9A6}

% ---- fixed font sizes, so layouts don't break in 11pt/12pt report classes ----
\newcommand{\uDiagScript}{\fontsize{7pt}{8.5pt}\selectfont}
\newcommand{\uDiagSmall}{\fontsize{9pt}{11pt}\selectfont}
\newcommand{\uDiagNormal}{\fontsize{10pt}{12pt}\selectfont}

\tikzset{
  % ---- generic ---------------------------------------------------------------
  every picture/.append style={line join=round,
     execute at begin picture={\let\scriptsize\uDiagScript
       \let\small\uDiagSmall\let\normalsize\uDiagNormal\uDiagSmall}},
  every node/.append style={execute at begin node={\hyphenpenalty=10000\exhyphenpenalty=10000}},
  uline/.style={draw=uLine, line width=0.6pt},
  lbl/.style={font=\scriptsize, text=uMuted, fill=white, inner sep=1pt},
  note/.style={draw=uMuted, fill=yellow!8, font=\scriptsize, align=left,
               text width=#1, inner sep=4pt, rounded corners=1pt},
  note/.default=4.6cm,
  % ---- use-case --------------------------------------------------------------
  usecase/.style={ellipse, draw=uLine, fill=white, line width=0.6pt,
                  minimum width=4.4cm, minimum height=0.95cm, align=center,
                  font=\scriptsize, inner sep=1pt, text width=3.6cm},
  sysboundary/.style={draw=uLine, line width=0.9pt, rounded corners=3pt,
                      fill=uIngest!25},
  assoc/.style={draw=uLine, line width=0.5pt},
  stereo/.style={draw=uLine, line width=0.5pt, dashed,
                 -{Stealth[length=2mm]}},
  extsys/.style={draw=uLine, fill=uGraph!60, rounded corners=2pt,
                 minimum width=2.3cm, minimum height=0.8cm, font=\scriptsize,
                 align=center},
  % ---- ERD / class boxes -------------------------------------------------------
  table/.style={rectangle split, rectangle split parts=2, draw=uLine,
                line width=0.6pt, rectangle split part fill={#1,white},
                text width=4.9cm, align=left, inner sep=3pt,
                font=\scriptsize\ttfamily},
  table/.default=uGraph,
  klass/.style={rectangle split, rectangle split parts=2, draw=uLine,
                line width=0.6pt, rectangle split part fill={#1,white},
                text width=3.7cm, align=center, inner sep=3pt,
                font=\scriptsize},
  klass/.default=white,
  rel/.style={draw=uLine, line width=0.55pt},
  card/.style={font=\scriptsize\bfseries, text=uInk, inner sep=1.5pt},
  generalise/.style={draw=uLine, line width=0.6pt,
                     -{Triangle[open, length=3mm, width=3mm]}},
  directed/.style={draw=uLine, line width=0.55pt, -{Stealth[length=2.2mm]}},
  % ---- flow ----------------------------------------------------------------------
  decision/.style={diamond, aspect=1.9, draw=uLine, fill=white,
                   line width=0.6pt, font=\scriptsize, align=center,
                   inner sep=1pt, text width=2.1cm},
  step/.style={rectangle, rounded corners=2pt, draw=uLine, line width=0.6pt,
               font=\scriptsize, align=center, minimum height=0.8cm,
               text width=#1, inner sep=3pt},
  step/.default=2.6cm,
  flow/.style={draw=uLine, line width=0.6pt, -{Stealth[length=2.2mm]}},
}

% ---- stick-figure actor:  \actor{name}{(x,y)}{Label} ---------------------------
\newcommand{\actor}[3]{%
  \begin{scope}[shift={#2}, line width=0.8pt, draw=uInk]
    \draw (0,0.62) circle (0.17);
    \draw (0,0.45) -- (0,-0.05);
    \draw (-0.32,0.28) -- (0.32,0.28);
    \draw (0,-0.05) -- (-0.24,-0.45);
    \draw (0,-0.05) -- (0.24,-0.45);
    \node[font=\scriptsize\bfseries, align=center, anchor=north]
         (#1-label) at (0,-0.5) {#3};
    \coordinate (#1) at (0.4,0.2);
  \end{scope}}

% Plain TikZ only (no tikz-uml dependency), so it compiles on any TeX Live.
% =============================================================================
\usepackage{tikz}
\usetikzlibrary{arrows.meta,positioning,calc,fit,backgrounds,
                shapes.geometric,shapes.multipart,decorations.pathreplacing}

% ---- palette (prints legibly in greyscale) ----------------------------------
\definecolor{uInk}{HTML}{1F2933}
\definecolor{uMuted}{HTML}{5B6670}
\definecolor{uLine}{HTML}{3E4C59}
\definecolor{uIngest}{HTML}{D9E8F5}   % event reliability / version store
\definecolor{uGraph}{HTML}{DDEFDF}    % evidence graph / policy / traceability
\definecolor{uVerify}{HTML}{F6E6CF}   % verification / findings
\definecolor{uIntent}{HTML}{E7E1F3}   % intent-side artifacts
\definecolor{uOk}{HTML}{2F7D4E}
\definecolor{uWarn}{HTML}{B7791F}
\definecolor{uBad}{HTML}{B4372F}
\definecolor{uSpan}{HTML}{FCE9A6}

% ---- fixed font sizes, so layouts don't break in 11pt/12pt report classes ----
\newcommand{\uDiagScript}{\fontsize{7pt}{8.5pt}\selectfont}
\newcommand{\uDiagSmall}{\fontsize{9pt}{11pt}\selectfont}
\newcommand{\uDiagNormal}{\fontsize{10pt}{12pt}\selectfont}

\tikzset{
  % ---- generic ---------------------------------------------------------------
  every picture/.append style={line join=round,
     execute at begin picture={\let\scriptsize\uDiagScript
       \let\small\uDiagSmall\let\normalsize\uDiagNormal\uDiagSmall}},
  every node/.append style={execute at begin node={\hyphenpenalty=10000\exhyphenpenalty=10000}},
  uline/.style={draw=uLine, line width=0.6pt},
  lbl/.style={font=\scriptsize, text=uMuted, fill=white, inner sep=1pt},
  note/.style={draw=uMuted, fill=yellow!8, font=\scriptsize, align=left,
               text width=#1, inner sep=4pt, rounded corners=1pt},
  note/.default=4.6cm,
  % ---- use-case --------------------------------------------------------------
  usecase/.style={ellipse, draw=uLine, fill=white, line width=0.6pt,
                  minimum width=4.4cm, minimum height=0.95cm, align=center,
                  font=\scriptsize, inner sep=1pt, text width=3.6cm},
  sysboundary/.style={draw=uLine, line width=0.9pt, rounded corners=3pt,
                      fill=uIngest!25},
  assoc/.style={draw=uLine, line width=0.5pt},
  stereo/.style={draw=uLine, line width=0.5pt, dashed,
                 -{Stealth[length=2mm]}},
  extsys/.style={draw=uLine, fill=uGraph!60, rounded corners=2pt,
                 minimum width=2.3cm, minimum height=0.8cm, font=\scriptsize,
                 align=center},
  % ---- ERD / class boxes -------------------------------------------------------
  table/.style={rectangle split, rectangle split parts=2, draw=uLine,
                line width=0.6pt, rectangle split part fill={#1,white},
                text width=4.9cm, align=left, inner sep=3pt,
                font=\scriptsize\ttfamily},
  table/.default=uGraph,
  klass/.style={rectangle split, rectangle split parts=2, draw=uLine,
                line width=0.6pt, rectangle split part fill={#1,white},
                text width=3.7cm, align=center, inner sep=3pt,
                font=\scriptsize},
  klass/.default=white,
  rel/.style={draw=uLine, line width=0.55pt},
  card/.style={font=\scriptsize\bfseries, text=uInk, inner sep=1.5pt},
  generalise/.style={draw=uLine, line width=0.6pt,
                     -{Triangle[open, length=3mm, width=3mm]}},
  directed/.style={draw=uLine, line width=0.55pt, -{Stealth[length=2.2mm]}},
  % ---- flow ----------------------------------------------------------------------
  decision/.style={diamond, aspect=1.9, draw=uLine, fill=white,
                   line width=0.6pt, font=\scriptsize, align=center,
                   inner sep=1pt, text width=2.1cm},
  step/.style={rectangle, rounded corners=2pt, draw=uLine, line width=0.6pt,
               font=\scriptsize, align=center, minimum height=0.8cm,
               text width=#1, inner sep=3pt},
  step/.default=2.6cm,
  flow/.style={draw=uLine, line width=0.6pt, -{Stealth[length=2.2mm]}},
}

% ---- stick-figure actor:  \actor{name}{(x,y)}{Label} ---------------------------
\newcommand{\actor}[3]{%
  \begin{scope}[shift={#2}, line width=0.8pt, draw=uInk]
    \draw (0,0.62) circle (0.17);
    \draw (0,0.45) -- (0,-0.05);
    \draw (-0.32,0.28) -- (0.32,0.28);
    \draw (0,-0.05) -- (-0.24,-0.45);
    \draw (0,-0.05) -- (0.24,-0.45);
    \node[font=\scriptsize\bfseries, align=center, anchor=north]
         (#1-label) at (0,-0.5) {#3};
    \coordinate (#1) at (0.4,0.2);
  \end{scope}}

% and this folder copied to  <report root>/diagrams/
% =============================================================================

% ---------------------------------------------------------------- ch4 --------
\begin{figure}[htbp]
  \centering
  \resizebox{\linewidth}{!}{% =============================================================================
% UNIFY -- Use-case diagram (owner: Mishal Fatima)
% Sources: proposal §1.1 (roles), §15.1 (FR-04, FR-11, FR-14, FR-17..FR-21),
%          §7.3 (designer-approved structured constraints),
%          §23.1 (PM approves pre-filled configuration; link confirmation goes
%          to code owner/assignee, not the PM), §26.2 (eligibility score).
% =============================================================================
\begin{tikzpicture}[x=1cm,y=1cm]

  % ---- system boundary --------------------------------------------------------
  \draw[sysboundary] (-1.0,-2.25) rectangle (11.2,10.6);
  \node[font=\small\bfseries, anchor=north west] at (-0.85,10.5)
       {UNIFY -- cross-tool change assurance};

  % ---- left column: human-facing use cases -------------------------------------
  \node[usecase] (sel)     at (1.9, 9.45) {Select change for verification\\(repo, base/head, linked artifacts)};
  \node[usecase] (evid)    at (1.9, 8.25) {View criterion-level evidence case};
  \node[usecase] (export)  at (1.9, 7.05) {Export audit evidence};
  \node[usecase] (link)    at (1.9, 5.85) {Confirm / reject cross-tool link};
  \node[usecase] (resolve) at (1.9, 4.6)  {Acknowledge / resolve / suppress / reopen finding};
  \node[usecase] (p4)      at (1.9, 3.35) {Review requirement--code divergence (P4)};
  \node[usecase] (notify)  at (1.9, 2.15) {Receive role-routed notification};
  \node[usecase] (constr)  at (1.9, 0.95) {Approve structured design constraint};
  \node[usecase] (config)  at (1.9,-0.25) {Approve pre-filled workflow configuration};
  \node[usecase] (connect) at (1.9,-1.45) {Connect repository and view eligibility score};

  % ---- right column -------------------------------------------------------------
  \node[usecase] (passes)  at (8.5, 8.25) {Inspect bounded verification passes};
  \node[usecase] (ingest)  at (8.5, 6.15) {Ingest and reconcile tool events};
  \node[usecase] (dispute) at (8.5, 4.6)  {Dispute verdict\\(logged as labelled correction)};

  % ---- include / extend -----------------------------------------------------------
  \draw[stereo] (passes) -- node[lbl, above=2pt]{\guillemotleft extend\guillemotright} (evid);
  \draw[stereo] (dispute) -- node[lbl, above=2pt]{\guillemotleft extend\guillemotright} (resolve);

  % ---- human actors (left) --------------------------------------------------------
  \actor{rev}{(-3.4,8.0)}{Code reviewer /\\technical lead\\(primary)}
  \actor{dev}{(-3.4,4.55)}{Developer}
  \actor{des}{(-3.4,1.35)}{Designer}
  \actor{pm}{(-3.4,-1.1)}{Product owner /\\PM}

  \foreach \u in {sel,evid,export,link,resolve} \draw[assoc] (rev) -- (\u.west);
  \foreach \u in {evid,link,resolve,p4,notify}  \draw[assoc] (dev) -- (\u.west);
  \foreach \u in {resolve,notify,constr}         \draw[assoc] (des) -- (\u.west);
  \foreach \u in {resolve,p4,notify,config,connect} \draw[assoc] (pm) -- (\u.west);

  % ---- external systems (right) ---------------------------------------------------
  \node[extsys] (gh)    at (13.6, 7.55) {\guillemotleft system\guillemotright\\GitHub};
  \node[extsys] (jira)  at (13.6, 6.15) {\guillemotleft system\guillemotright\\Jira};
  \node[extsys] (figma) at (13.6, 4.75) {\guillemotleft system\guillemotright\\Figma};
  \node[extsys] (slack) at (13.6, 2.15)  {\guillemotleft system\guillemotright\\Slack (delivery only)};

  \foreach \s in {gh,jira,figma} \draw[assoc] (\s.west) -- (ingest);
  \draw[assoc] (slack.west) -- (notify);

  % ---- note ----------------------------------------------------------------------
  \node[note=4.6cm, anchor=north west] at (6.0,1.3)
    {Link confirmation is routed to the code owner/assignee, not the PM (\S23.1).
     Configuration is proposed by the system from observed transitions; the PM
     only approves it. Slack never feeds reasoning.};
\end{tikzpicture}
}
  \caption{Use-case diagram for UNIFY's four stakeholder roles and the three
           source tools. Link confirmation is routed to the code owner or
           assignee rather than the PM, and configuration is system-proposed
           and PM-approved.}
  \label{fig:usecase}
\end{figure}

% ---------------------------------------------------------------- ch5 --------
\begin{landscape}
\begin{figure}[p]
  \centering
  \resizebox{\linewidth}{!}{% =============================================================================
% UNIFY -- Domain / class model of the evidence graph (owner: Mishal Fatima)
% Sources: proposal §9.1 (node types, edge types, edge attributes), §9.2
%          (eligibility of edges), §7.3 (structured design constraints),
%          §22.1-22.2 (Requirement, SATISFIES + CodeAnchor), §26.1 (symbol card).
% Every association drawn here is persisted as a row of `edges` (see ERD).
% =============================================================================
\begin{tikzpicture}[x=1cm,y=1cm]
  \tikzset{klass/.append style={align=left, text width=3.45cm}}
  \newcommand{\cname}[2][]{\makebox[\linewidth]{\textbf{#2}}%
     \ifx\relax#1\relax\else\\[-1pt]\makebox[\linewidth]{\textit{\scriptsize#1}}\fi}

  % ---------------- abstract superclass + person -------------------------------
  \node[klass=uGraph, anchor=center] (art) at (-2.6,0) {%
    \cname[\guillemotleft abstract\guillemotright]{Artifact}
    \nodepart{two}
    id : UUID\\ externalId : String\\ version : String\\ status : String\\
    owner : String\\ lastModified : Time\\ validFrom, validTo : Time\\ raw : JSON};

  \node[klass=white, anchor=north] (person) at (-2.6,-4.6) {%
    \cname{Person}
    \nodepart{two}
    handle : String\\ displayName : String\\ toolIdentities : Map};

  \draw[directed] (person.north) -- (art.south)
    node[lbl, pos=0.5, align=center]{OWNS / ASSIGNED\_TO\\/ WATCHES};
  \node[card, anchor=north west] at ([xshift=2pt]art.south) {*};
  \node[card, anchor=south west] at ([xshift=2pt]person.north) {*};

  \node[note=4.4cm, anchor=south] (edgenote) at (-2.6,2.6) {%
    \textbf{Every association is an \texttt{Edge}} carrying
    \emph{sourceClass} (observed | derived | model\_candidate | unresolved),
    confidence, resolver, createdAt and a validity interval.
    Only deterministic or confirmed edges may produce a decisive finding (\S9.2).};

  % ---------------- row A: intent side (anchored on their south edge) ----------
  \node[klass=uIntent, anchor=south] (da) at (2.6,3) {%
    \cname{DesignArtifact}\nodepart{two} fileKey : String\\ nodeId : String\\ versionId : String};
  \node[klass=uIntent, anchor=south] (dc) at (7.2,3) {%
    \cname{DesignConstraint}\nodepart{two} kind : ConstraintKind\\ expected : String\\ approvedBy : Person};
  \node[klass=uIntent, anchor=south] (tk) at (11.8,3) {%
    \cname{Ticket}\nodepart{two} key : String\\ workflowState : String};
  \node[klass=uIntent, anchor=south] (cr) at (16.4,3) {%
    \cname{Criterion}\nodepart{two} text : String\\ atomicIndex : Int};
  \node[klass=white, anchor=south] (msg) at (21.0,3) {%
    \cname{Message}\nodepart{two} channel : String\\ ts : Time};

  % ---------------- row B: change / evidence side (anchored on north edge) -----
  \node[klass=uIngest, anchor=north] (pr) at (2.6,-3) {%
    \cname{PullRequest}\nodepart{two} number : Int\\ baseSha, headSha : String\\ mergedAt : Time [0..1]};
  \node[klass=uIngest, anchor=north] (cm) at (7.2,-3) {%
    \cname{Commit}\nodepart{two} sha : String\\ message : String\\ authoredAt : Time};
  \node[klass=uIngest, anchor=north] (tc) at (11.8,-3) {%
    \cname{TestCase}\nodepart{two} path : String\\ name : String\\ lastResult : Result [0..1]};
  \node[klass=uIngest, anchor=north] (cs) at (16.4,-3) {%
    \cname[symbol card fields, \S26.1]{CodeSymbol}\nodepart{two}
    fqName : String\\ signature : String\\ filePath : String\\ docstring : String [0..1]\\
    blobHash : String};
  \node[klass=uIntent, anchor=north] (rq) at (22.4,-3) {%
    \cname{Requirement}\nodepart{two} reqKey : String\\ kind : FR | NFR\\ text : String\\ source : String};

  % ---------------- generalisation tree ----------------------------------------
  \coordinate (trunkTop) at (-0.25,5.5);
  \coordinate (trunkBot) at (-0.25,-7.1);
  \draw[uline] (trunkTop) -- (trunkBot);
  \draw[uline] (trunkTop) -- (21.0,5.5);
  \draw[uline] (trunkBot) -- (22.4,-7.1);
  \foreach \c in {da,dc,tk,cr,msg} \draw[uline] (\c.north) -- (\c.north |- trunkTop);
  \foreach \c in {pr,cm,tc,cs,rq}  \draw[uline] (\c.south) -- (\c.south |- trunkBot);
  \draw[generalise] (-0.25,0) -- (art.east);

  % ---------------- associations within row A ----------------------------------
  \draw[rel] ([yshift=0.45cm]dc.south west) -- ([yshift=0.45cm]dc.south west -| da.east)
    node[lbl, pos=0.5, above=1pt]{PART\_OF}
    node[card, pos=0.05, anchor=north east]{*}
    node[card, pos=0.95, anchor=north west]{1};
  \draw[rel] ([yshift=0.45cm]cr.south west) -- ([yshift=0.45cm]cr.south west -| tk.east)
    node[lbl, pos=0.5, above=1pt]{PART\_OF}
    node[card, pos=0.05, anchor=north east]{1..*}
    node[card, pos=0.95, anchor=north west]{1};

  % CONSTRAINS: design constraint -> criterion (routed below row A)
  \draw[directed] (dc.south) |- (15.3,1.7) -- (15.3,3)
    node[lbl, pos=0.5, anchor=south]{CONSTRAINS};

  % ---------------- associations within row B ----------------------------------
  \draw[directed] ([yshift=-0.45cm]cm.north west) -- ([yshift=-0.45cm]cm.north west -| pr.east)
    node[lbl, pos=0.5, below=1pt]{PART\_OF}
    node[card, pos=0.06, anchor=south east]{1..*}
    node[card, pos=0.94, anchor=south west]{1};
  \draw[directed] ([yshift=-0.45cm]tc.north east) -- ([yshift=-0.45cm]tc.north east -| cs.west)
    node[lbl, pos=0.5, below=1pt]{COVERS}
    node[card, pos=0.06, anchor=south west]{*}
    node[card, pos=0.94, anchor=south east]{*};

  % TOUCHES_COMPONENT: self-association on PullRequest (top-left corner)
  \draw[uline] ([xshift=0.35cm]pr.north west) -- ++(0,0.55) -| ([xshift=-0.45cm,yshift=-0.5cm]pr.north west)
     -- ([yshift=-0.5cm]pr.north west);
  \node[lbl, anchor=south, align=center] at ([xshift=-0.05cm,yshift=0.6cm]pr.north west)
       {TOUCHES\_COMPONENT};

  % ---------------- associations across rows -----------------------------------
  \draw[directed] (pr.north) -- (da.south)
    node[lbl, pos=0.5, align=center]{IMPLEMENTS\\(design version)};
  \draw[directed] ([xshift=0.9cm]pr.north) -- ([xshift=-0.8cm]tk.south)
    node[lbl, pos=0.62, sloped, above=1pt]{TRACKS};
  \draw[directed] ([xshift=0.6cm]tc.north) -- ([xshift=-0.4cm]cr.south)
    node[lbl, pos=0.45, sloped, above=1pt]{COVERS};
  \draw[directed, dashed] ([xshift=0.8cm]cs.north) -- ([xshift=0.8cm]cr.south)
    node[lbl, pos=0.5, align=right, anchor=east, xshift=-2pt]{SATISFIES\_\\CANDIDATE\\\textit{(never decisive}\\\textit{on its own)}};

  % SATISFIES with association class CodeAnchor
  \draw[directed] ([yshift=-0.6cm]cs.north east) -- ([yshift=-0.6cm]cs.north east -| rq.west)
    node[lbl, pos=0.5, below=2pt]{SATISFIES}
    node[card, pos=0.04, anchor=south west]{*}
    node[card, pos=0.96, anchor=south east]{*};
  \node[klass=uGraph, anchor=south] (anchor) at (19.45,-1.75) {%
    \cname[\guillemotleft association class\guillemotright]{CodeAnchor}
    \nodepart{two}
    repo, filePath : String\\ symbolFqn : String\\ startLine, endLine : Int\\
    blobHash, commitSha : String\\ method : AnchorMethod\\ status : ACTIVE | BROKEN\\
    confidence : Real};
  \coordinate (satline) at ([yshift=-0.6cm]cs.north east);
  \draw[uline, dashed] (anchor.south) -- (anchor.south |- satline);

  % ---------------- notes ------------------------------------------------------
  \node[note=3.3cm, anchor=west] (msgnote) at (23.1,4.0)
    {Any \texttt{Artifact} may be \textbf{DISCUSSED\_IN} a \texttt{Message}.
     Messages are a role/watch signal only, never a rule-evaluation target (\S9.1).};
  \draw[uline, dashed] (msgnote.west) -- (msg.east);

  \node[note=3.3cm, anchor=north west] (nfrnote) at (24.5,-3)
    {\textbf{FR} anchors to symbol + line range, commit-pinned.
     \textbf{NFR} is \emph{not} line-anchored: it links to its enforcement point
     and verifying \texttt{TestCase} (\S22.4).};
  \draw[uline, dashed] (nfrnote.west) -- (rq.east);

  % ---------------- legend -----------------------------------------------------
  \node[note=4.4cm, anchor=north] at (-2.6,-8.1) {%
    \tikz\fill[uIntent,draw=uLine] (0,0) rectangle (0.35,0.22);\ linked intent (what was asked for)\\
    \tikz\fill[uIngest,draw=uLine] (0,0) rectangle (0.35,0.22);\ actual change and test evidence\\
    \tikz\fill[uGraph,draw=uLine] (0,0) rectangle (0.35,0.22);\ graph structure / traceability\\[2pt]
    Solid arrow: observed or derived edge. Dashed arrow: retrieval candidate.};
\end{tikzpicture}
}
  \caption{Domain model of the evidence graph. Every association is persisted
           as a typed, provenance-carrying edge; \texttt{CodeAnchor} is the
           association class on \texttt{SATISFIES}.}
  \label{fig:domain-model}
\end{figure}
\end{landscape}

\begin{landscape}
\begin{figure}[p]
  \centering
  \resizebox{\linewidth}{!}{% =============================================================================
% UNIFY -- Entity-relationship diagram, PostgreSQL MVP (owner: Mishal Fatima)
% Sources: proposal §8.3 (event envelope), §8.4 (checkpoints, SOURCE_STALE),
%          §9.1/§9.9 (nodes/edges/artifact_snapshots + indexes), §9.3 (rule
%          schema), §11.5-11.8 (packet, verdict schema, gate), §15.1 FR-03,
%          FR-15, FR-18 (immutable evidence, idempotent findings, audit),
%          §22.1-22.2 (CodeAnchor, re-anchoring), §27.1 (packet manifest).
% Design decision: Requirement, Criterion, Ticket, ... are rows of `nodes`
% (type discriminator + jsonb payload), exactly as §9.9 commits to. Only the
% CodeAnchor gets its own table, because re-anchoring history and
% anchor_survival_rate need one row per commit-pinned anchor.
% =============================================================================
\begin{tikzpicture}[x=1cm,y=1cm]
  \newcommand{\pk}{\textbf{PK}\ }
  \newcommand{\fk}{\textit{FK}\ }
  \newcommand{\ix}{\phantom{PK}\ }
  \newcommand{\rarr}{$\rightarrow$}
  \newcommand{\hdr}[2]{\textbf{\normalsize #1}\ \textrm{\scriptsize\textit{#2}}}

  % ============================ ROW 1 =========================================
  \node[table=uIngest, anchor=north] (events) at (-13,0) {%
    \hdr{events}{\guillemotleft append-only\guillemotright}
    \nodepart{two}
    \pk event\_id : text\\
    \ix vendor\_delivery\_id : text\\
    \ix workspace\_id : uuid\\
    \ix source\_tool : enum\\
    \ix entity\_type, entity\_id : text\\
    \ix version : text\\
    \ix action : text\\
    \ix state\_before, state\_after : text\\
    \ix actor : text\\
    \ix occurred\_at, received\_at : tstz\\
    \ix references : jsonb\\
    \ix raw\_hash : text\\
    \ix schema\_version : int\\[2pt]
    \textrm{\textit{UNIQUE (source\_tool, vendor\_delivery\_id)}}};

  \node[table=uIngest, anchor=north] (snaps) at (-6.5,0) {%
    \hdr{artifact\_snapshots}{\guillemotleft immutable\guillemotright}
    \nodepart{two}
    \pk id : uuid\\
    \fk node\_id \rarr{} nodes\\
    \fk source\_event\_id \rarr{} events\\
    \ix version : text\\
    \ix content\_hash : text\\
    \ix captured\_at : tstz\\
    \ix content : jsonb\\[2pt]
    \textrm{\textit{UNIQUE (node\_id, version)}}};

  \node[table=uGraph, anchor=north] (nodes) at (0,0) {%
    \hdr{nodes}{}
    \nodepart{two}
    \pk id : uuid\\
    \ix workspace\_id : uuid\\
    \ix type : node\_type (enum)\\
    \ix external\_id : text\\
    \ix version : text\\
    \ix status, owner : text\\
    \ix last\_modified : tstz\\
    \ix valid\_from, valid\_to : tstz\\
    \ix payload : jsonb\\[2pt]
    \textrm{\textit{UNIQUE (workspace\_id, type, external\_id)}}};

  \node[table=uGraph, anchor=north] (edges) at (6.5,0) {%
    \hdr{edges}{}
    \nodepart{two}
    \pk id : uuid\\
    \fk from\_id \rarr{} nodes\\
    \fk to\_id \rarr{} nodes\\
    \ix type : edge\_type (enum)\\
    \ix source\_class : enum\\
    \ix confidence : real\\
    \ix resolver : text\\
    \ix link\_state : enum\\
    \ix created\_at : tstz\\
    \ix valid\_from, valid\_to : tstz\\
    \ix metadata : jsonb\\[2pt]
    \textrm{\textit{INDEX (from\_id, type), (to\_id, type)}}\\
    \textrm{\textit{UNIQUE (from\_id, to\_id, type)}}\\
    \textrm{\textit{\hphantom{UNIQUE }\textcolor{uMuted}{WHERE link\_state $\neq$ 'revoked'}}}};

  \node[table=uGraph, anchor=north] (anchors) at (13,0) {%
    \hdr{code\_anchors}{\guillemotleft commit-pinned\guillemotright}
    \nodepart{two}
    \pk id : uuid\\
    \fk edge\_id \rarr{} edges (SATISFIES)\\
    \fk prev\_anchor\_id \rarr{} code\_anchors\\
    \ix repo, file\_path : text\\
    \ix symbol\_fqn : text\\
    \ix start\_line, end\_line : int\\
    \ix blob\_hash : text\\
    \ix commit\_sha : text\\
    \ix method : anchor\_method (enum)\\
    \ix status : anchor\_status (enum)\\
    \ix confidence : real\\
    \ix provenance : jsonb};

  % ============================ ROW 2 =========================================
  \node[table=uIngest, anchor=north] (ckpt) at (-13,-7.0) {%
    \hdr{checkpoints}{}
    \nodepart{two}
    \pk workspace\_id, source\_tool,\\
    \ix resource\_id\\
    \fk last\_event\_id \rarr{} events\\
    \ix last\_state\_at : tstz\\
    \ix last\_success\_at : tstz\\
    \ix source\_stale : bool};

  \node[table=uVerify, anchor=north] (fe) at (-6.5,-7.0) {%
    \hdr{finding\_evidence}{}
    \nodepart{two}
    \pk (finding\_id, snapshot\_id)\\
    \fk finding\_id \rarr{} findings\\
    \fk snapshot\_id \rarr{} artifact\_snapshots\\
    \ix role : text};

  \node[table=uVerify, anchor=north] (findings) at (0,-7.0) {%
    \hdr{findings}{}
    \nodepart{two}
    \pk id : uuid\\
    \fk subject\_node\_id \rarr{} nodes\\
    \fk rule\_id, rule\_version \rarr{} rules\\
    \fk verdict\_id \rarr{} verdicts\\
    \fk superseded\_by\_id \rarr{} findings\\
    \ix kind : enum (rule | verdict)\\
    \ix idempotency\_key : text\\
    \ix annotation, severity : text\\
    \ix state : finding\_state (enum)\\
    \ix created\_at : tstz\\[2pt]
    \textrm{\textit{UNIQUE (idempotency\_key)}}\\
    \textrm{\textit{CHECK exactly one of rule / verdict}}\\
    \textrm{\textit{CHECK state = superseded}}\\
    \textrm{\textit{\hphantom{CHECK }$\iff$ superseded\_by\_id set}}};

  \node[table=uVerify, anchor=north] (verdicts) at (6.5,-7.0) {%
    \hdr{verdicts}{}
    \nodepart{two}
    \pk id : uuid\\
    \fk packet\_id \rarr{} context\_packets\\
    \ix status : verdict\_status (enum)\\
    \ix evidence : jsonb (cited spans)\\
    \ix reason, missing\_evidence : text\\
    \ix divergence\_tag : enum\\
    \ix freshness\_state : enum\\
    \ix review\_required : bool\\
    \ix reason\_code : text\\
    \ix pass\_count : int, stop\_reason\\
    \ix gate\_outcome : enum\\
    \ix model\_and\_run : jsonb};

  \node[table=uVerify, anchor=north] (packets) at (13,-7.0) {%
    \hdr{context\_packets}{\guillemotleft immutable\guillemotright}
    \nodepart{two}
    \pk id : uuid\\
    \fk criterion\_node\_id \rarr{} nodes\\
    \ix base\_commit, head\_commit\\
    \ix token\_budget : int\\
    \ix slot\_budgets : jsonb\\
    \ix manifest : jsonb\\
    \ix prompt\_hash : text\\
    \ix packet\_hash : text};

  % ============================ ROW 3 =========================================
  \node[table=uGraph, anchor=north] (rules) at (-6.5,-13.9) {%
    \hdr{rules}{\guillemotleft versioned\guillemotright}
    \nodepart{two}
    \pk id, version\\
    \ix edge\_type : edge\_type\\
    \ix trigger, guards : jsonb\\
    \ix condition : text\\
    \ix grace\_period\_hours : int\\
    \ix produces, severity : text\\
    \ix notify : jsonb\\
    \ix resolution\_options : text[]\\
    \ix active : bool};

  \node[table=uVerify, anchor=north] (audit) at (0,-13.9) {%
    \hdr{audit\_log}{\guillemotleft append-only\guillemotright}
    \nodepart{two}
    \pk id : uuid\\
    \fk finding\_id \rarr{} findings\\
    \fk actor\_id \rarr{} nodes (Person)\\
    \ix action : enum\\
    \ix reason : text\\
    \ix evidence\_snapshot : jsonb\\
    \ix at : tstz};

  \node[table=uVerify, anchor=north] (notif) at (6.5,-13.9) {%
    \hdr{notifications}{}
    \nodepart{two}
    \pk id : uuid\\
    \fk finding\_id \rarr{} findings\\
    \fk recipient\_id \rarr{} nodes (Person)\\
    \ix channel : text\\
    \ix delivery\_status : enum\\
    \ix delivered\_at : tstz};

  % ============================ RELATIONSHIPS ==================================
  % helper: horizontal relationship at a given depth below the row top
  \newcommand{\hrel}[6]{% from, to, depth, card-from, card-to, label
    \draw[rel] ([yshift=-#3]#1.north east) -- ([yshift=-#3]#1.north east -| #2.west)
      node[card, pos=0.0, anchor=south west]{#4}
      node[card, pos=1.0, anchor=south east]{#5}
      node[lbl, pos=0.5, below=1pt]{#6};}

  \hrel{events}{snaps}{1.2cm}{1}{0..*}{source}
  \hrel{snaps}{nodes}{1.2cm}{1..*}{1}{of}
  \hrel{nodes}{edges}{1.0cm}{1}{0..*}{from}
  \hrel{nodes}{edges}{2.0cm}{1}{0..*}{to}
  \hrel{edges}{anchors}{1.2cm}{1}{0..*}{SATISFIES}
  \hrel{fe}{findings}{1.0cm}{1..*}{1}{cites}
  \hrel{findings}{verdicts}{1.2cm}{0..1}{0..1}{from}
  \hrel{verdicts}{packets}{1.2cm}{1..*}{1}{judges}

  % self-reference: superseded_by_id (re-verification chain, finding lifecycle)
  \draw[rel] ([yshift=-2.6cm]findings.north west) -- ++(-0.7,0) |- ([yshift=-1.9cm]findings.north west);
  \node[card, anchor=east, align=right, font=\scriptsize] at ([xshift=-0.75cm,yshift=-2.22cm]findings.north west)
       {0..1\\[-1pt]\textmd{\color{uMuted}newer}};

  % self-reference: prev_anchor_id (re-anchoring history)
  \draw[rel] ([yshift=-1.4cm]anchors.north east) -- ++(0.6,0) |- ([yshift=-0.75cm]anchors.north east);
  \node[card, anchor=west, align=left, font=\scriptsize] at ([xshift=0.65cm,yshift=-1.07cm]anchors.north east)
       {0..1\\[-1pt]\textmd{\color{uMuted}prev}};

  % vertical relationships
  \draw[rel] (ckpt.north) -- (events.south)
    node[card, pos=0.04, anchor=south west]{0..*}
    node[card, pos=0.96, anchor=north west]{0..1};
  \draw[rel] (fe.north) -- (snaps.south)
    node[card, pos=0.06, anchor=south west]{0..*}
    node[card, pos=0.95, anchor=north west]{1};
  \draw[rel] (findings.north) -- (nodes.south)
    node[card, pos=0.06, anchor=south west]{0..*}
    node[card, pos=0.95, anchor=north west]{1}
    node[lbl, pos=0.5, right=1pt]{subject};
  \draw[rel] (audit.north) -- (findings.south)
    node[card, pos=0.1, anchor=south west]{0..*}
    node[card, pos=0.9, anchor=north west]{1};

  % findings -> rules (orthogonal, below finding_evidence)
  \draw[rel] ([yshift=0.5cm]findings.south west) -| (rules.north)
    node[card, pos=0.03, anchor=south east]{0..*}
    node[card, pos=0.97, anchor=south west]{0..1}
    node[lbl, pos=0.3, above=1pt]{fired by};

  % findings -> notifications
  \draw[rel] ([xshift=1.6cm]findings.south) |- ++(4.9,-0.75) -- (notif.north)
    node[card, pos=0.0, anchor=north west]{0..*};
  \node[card, anchor=south west] at ([xshift=1.6cm]findings.south) {1};

  % context_packets -> nodes (criterion), routed over the top row
  \draw[rel] ([yshift=-1.0cm]packets.north east) -- ++(2.3,0) |- (0.6,1.0) -- (0.6,0)
    node[card, pos=0.85, anchor=west]{1};
  \node[card, anchor=south west] at ([yshift=-1.0cm]packets.north east) {0..*};
  \node[lbl, anchor=south] at (8,1.0) {criterion (Criterion node)};

  % ============================ LEGEND ==========================================
  \node[note=5.4cm, anchor=north] (legend) at (-13,-12.3) {%
    \textbf{Layer (header colour)}\\[2pt]
    \tikz\fill[uIngest,draw=uLine] (0,0) rectangle (0.35,0.22);\ event reliability \& version store\\
    \tikz\fill[uGraph,draw=uLine] (0,0) rectangle (0.35,0.22);\ evidence graph, policy \& traceability\\
    \tikz\fill[uVerify,draw=uLine] (0,0) rectangle (0.35,0.22);\ verification \& findings\\[4pt]
    \textbf{Notation}: \texttt{PK} primary key, \texttt{\textit{FK}} foreign key,
    \texttt{1}/\texttt{0..1}/\texttt{0..*}/\texttt{1..*} cardinality at that end,
    \texttt{tstz} = timestamptz.\\[4pt]
    \textbf{Enums}\\
    \texttt{node\_type}: DesignArtifact, DesignConstraint, Ticket, Criterion,
      Requirement, PullRequest, Commit, CodeSymbol, TestCase, Person, Message\\
    \texttt{source\_class}: observed | derived | model\_candidate | unresolved\\
    \texttt{link\_state}: active | dismissed | revoked\\
    \texttt{anchor\_method}: created | symbol | hash | hunk\\
    \texttt{anchor\_status}: active | broken\\
    \texttt{verdict\_status}: SUPPORTED | CONTRADICTED | PARTIAL | INSUFFICIENT\_EVIDENCE\\
    \texttt{finding\_state}: open | acknowledged | resolved | disputed |
      suppressed | superseded};
\end{tikzpicture}
}
  \caption{Entity-relationship diagram of the PostgreSQL schema. Graph entities
           share one \texttt{nodes} table with a type discriminator; evidence
           tables are append-only or immutable.}
  \label{fig:erd}
\end{figure}
\end{landscape}

\begin{landscape}
\begin{figure}[p]
  \centering
  \resizebox{\linewidth}{!}{% =============================================================================
% UNIFY -- Requirement -> code anchor: commit-pinned anchoring, re-anchoring
% order and the broken state (owner: Mishal Fatima)
% Sources: proposal §22.1 (CodeAnchor fields), §22.2 (re-anchor order:
%          fully-qualified symbol -> content hash -> diff-hunk mapping -> broken;
%          anchor survival rate), §22.3 (P4), §22.4 (FR vs NFR granularity).
% The file names, symbols and line numbers are an ILLUSTRATIVE example only.
% =============================================================================
\begin{tikzpicture}[x=1cm,y=1cm]

  % ---- one commit panel ---------------------------------------------------------
  % #1 x-origin, #2 commit id, #3 event caption, #4 file name, #5 span top line
  % index (0-11), #6 span bottom line index, #7 span label, #8 tag colour,
  % #9 tag text. Set \badge to the decision-step number first (empty = none).
  \newcommand{\badge}{}%
  \newcommand{\panel}[9]{%
    \begin{scope}[shift={(#1,0)}]
      \node[circle, draw=uLine, fill=white, line width=0.7pt, minimum size=0.75cm,
            font=\scriptsize\bfseries] (#2) at (2.2,0) {#2};
      \node[font=\scriptsize, align=center, text width=4.3cm, anchor=south]
            at (2.2,0.45) {#3};
      % file strip
      \draw[uLine, fill=white, line width=0.6pt] (0,-0.75) rectangle (4.4,-3.95);
      \fill[uIngest] (0,-0.75) rectangle (4.4,-1.15);
      \draw[uLine, line width=0.6pt] (0,-1.15) -- (4.4,-1.15);
      \node[font=\scriptsize\ttfamily, anchor=west] at (0.08,-0.95) {#4};
      \ifnum#5>-1
        \fill[uSpan, draw=uWarn, line width=0.5pt]
          (0.05,{-1.3-0.22*#5+0.08}) rectangle (4.35,{-1.3-0.22*#6-0.08});
      \fi
      \foreach \i/\l in {0/2.6,1/3.4,2/1.9,3/3.0,4/2.4,5/3.6,6/2.8,7/1.6,8/3.1,9/2.2,10/2.9,11/1.4}
        \draw[uMuted!55, line width=1.6pt] (0.3,{-1.3-0.22*\i}) -- ++(\l,0);
      \ifnum#5>-1
        \node[font=\scriptsize\ttfamily, fill=uSpan, inner sep=1pt, anchor=east]
          at (4.3,{-1.3-0.22*(#5+#6)/2}) {#7};
      \else
        \node[font=\scriptsize, text=uBad, fill=white, inner sep=4pt, align=center,
              minimum width=4.3cm, minimum height=0.7cm]
          at (2.2,-2.5) {#7};
      \fi
      % status tag
      \node[draw=#8, fill=#8!12, text=#8!70!black, rounded corners=2pt,
            font=\scriptsize\bfseries, align=center, text width=4.1cm,
            anchor=north, inner sep=3pt] (#2tag) at (2.2,-4.15) {#9};
      \if\relax\detokenize\expandafter{\badge}\relax\else
        \node[circle, fill=uInk, text=white, font=\scriptsize\bfseries,
              inner sep=1pt, minimum size=0.42cm] at (4.4,-0.75) {\badge};
      \fi
    \end{scope}}

  % ---- the five commits ----------------------------------------------------------
  \renewcommand{\badge}{}
  \panel{0}{c1}{anchor created\\(SATISFIES confirmed)}
        {src/cart/validate.ts}{2}{6}{validateCart()\ \ L40--58}
        {uOk}{ACTIVE\\\textmd{method = created}}
  \renewcommand{\badge}{1}
  \panel{5.3}{c2}{12 lines inserted\\above the function}
        {src/cart/validate.ts}{5}{9}{validateCart()\ \ L52--70}
        {uOk}{ACTIVE\\\textmd{re-anchored by symbol name}}
  \renewcommand{\badge}{2}
  \panel{10.6}{c3}{function renamed,\\body unchanged}
        {src/cart/validate.ts}{5}{9}{checkCart()\ \ L52--70}
        {uOk}{ACTIVE\\\textmd{re-anchored by content hash}}
  \renewcommand{\badge}{3}
  \panel{15.9}{c4}{logic inlined into\\submitOrder() and edited}
        {src/cart/submit.ts}{6}{10}{submitOrder()\ \ L112--127}
        {uWarn}{ACTIVE, lower confidence\\\textmd{re-anchored by diff-hunk map}}
  \renewcommand{\badge}{4}
  \panel{21.2}{c5}{logic deleted}
        {src/cart/submit.ts}{-1}{-1}{anchored span no longer exists}
        {uBad}{BROKEN\\\textmd{shown to the user, never silently kept}}

  \foreach \a/\b in {c1/c2,c2/c3,c3/c4,c4/c5} \draw[flow] (\a) -- (\b);
  \draw[flow, dashed] (c5) -- ++(1.6,0) node[right, font=\scriptsize, text=uMuted]{time};

  % ---- requirement on the left ----------------------------------------------------
  \node[draw=uLine, fill=uIntent, line width=0.6pt, rounded corners=2pt,
        text width=3.2cm, align=left, font=\scriptsize, inner sep=4pt, anchor=east]
        (req) at (-1.9,-2.3)
        {\textbf{Requirement FR-07} (v3)\\[2pt]
         ``Checkout shall reject an empty cart.''\\[3pt]
         \textit{kind:} FR $\Rightarrow$ line-anchored};
  \draw[directed] (req.east) -- (0,-2.3)
    node[lbl, pos=0.45, above=2pt, align=center]{SATISFIES\\(pinned to c1)};

  % ---- anchor tuple callout ----------------------------------------------------------
  \node[note=3.6cm, anchor=north east] (tuple) at (-1.9,-3.9)
    {\textbf{CodeAnchor @ c1}\\
     repo: acme/storefront\\
     file: src/cart/validate.ts\\
     symbol: cart.validateCart\\
     lines: 40--58\\
     blob: 9f2c\ldots{} \quad commit: c1\\
     confidence: 1.0 (deterministic)\\
     provenance: confirmed by reviewer};

  % ---- re-anchoring decision ladder ----------------------------------------------------
  \begin{scope}[shift={(0,-8.0)}]
    \node[step=2.9cm, fill=uIngest] (new) at (0.9,0)
      {New commit touches the anchored file};
    \node[decision] (d1) at (5.4,0) {\textbf{1} Fully-qualified symbol still present?};
    \node[decision] (d2) at (10.7,0) {\textbf{2} Identical content hash found?};
    \node[decision] (d3) at (16.0,0) {\textbf{3} Diff hunk maps the old span?};
    \node[step=3.3cm, draw=uBad, fill=uBad!10, text=uBad!70!black] (brk) at (21.4,0)
      {\textbf{4} Mark anchor \textbf{BROKEN}, keep history, raise a finding};

    \draw[flow] (new) -- (d1);
    \draw[flow] (d1) -- node[lbl, above=2pt]{no} (d2);
    \draw[flow] (d2) -- node[lbl, above=2pt]{no} (d3);
    \draw[flow] (d3) -- node[lbl, above=2pt]{no} (brk);

    \node[step=3.2cm, draw=uOk, fill=uOk!10] (r1) at (5.4,-2.3)
      {Re-anchor to the symbol's new line range};
    \node[step=3.2cm, draw=uOk, fill=uOk!10] (r2) at (10.7,-2.3)
      {Re-anchor to the moved / renamed span};
    \node[step=3.2cm, draw=uWarn, fill=uWarn!10] (r3) at (16.0,-2.3)
      {Re-anchor to the mapped span; lower confidence};
    \foreach \d/\r in {d1/r1,d2/r2,d3/r3}
      \draw[flow] (\d) -- node[lbl, right=2pt]{yes} (\r);

    \node[note=3.6cm, anchor=north] at (21.4,-1.0)
      {Every outcome appends a new \texttt{code\_anchors} row (\texttt{prev\_anchor\_id}
       links the history), so a finding can always be re-read against the commit it cited.\\[3pt]
       \textbf{anchor\_survival\_rate} = anchors not broken $\div$ anchors tracked,
       reported per commit window (\S22.2).};

    \node[note=3.6cm, anchor=north east] at (-1.9,-0.4)
      {\textbf{NFRs are not line-anchored.}
       ``Respond within 2\,s'' links to its \emph{enforcement point}
       (timeout constant, queue config) and a \emph{verifying test};
       its metric is coverage, not span overlap (\S22.4).};
  \end{scope}
\end{tikzpicture}
}
  \caption{Commit-pinned requirement-to-code anchoring (illustrative example).
           Re-anchoring tries the fully-qualified symbol, then the content hash,
           then diff-hunk mapping; if all fail, the anchor is marked broken
           rather than left silently wrong.}
  \label{fig:anchor}
\end{figure}
\end{landscape}
